% Módulo 02e — Exegese: etapas, evidências, cálculos, lemas e estado da arte (R776)
\section{Leitura guiada do trabalho: por que cada escolha}
\textit{Etapas.} O trabalho encadeou-se assim, cada fase respondendo à anterior: núcleo de verificação (R747) porque alegação sem verificador é retórica; projeto lake com Mathlib (R754) porque compilador sem biblioteca não toca teoremas reais; teoremas de calibração (T1--T7, R755) porque pesos sem semântica são números órfãos; invariância, pertinência bilateral e singleton porque a revisão exigiu casos-limite; monotonicidade e top-\textit{k} (R761) porque decidir pelo máximo pedia justificativa; temperatura e overflow-free (R763) porque a literatura e o código os exigiam; pipeline de sincronia (R767) porque deriva silenciosa já mordeu uma vez.
\textit{Evidências.} Lean com Mathlib: o único par linguagem-biblioteca com kernel verificador e biblioteca matemática madura acessível neste ambiente \cite{moura2021,mathlib2020}. Bancada sintética: oráculo oculto para não-circularidade, com arbitrariedade declarada e sensibilidade em três mundos como contrapeso. Intervalos por aproximação normal com $n=200$ e pareamento por decisão: o desenho mais simples que ainda separa sinal de ruído. Fuzz em duas sementes: a forma mais barata de interrogar floats. Pinos e SHAs: identidade criptográfica de terceiros sem executar terceiros.
\textit{Cálculos.} Arrependimento mede perda contra o ótimo do oráculo; IC95 quantifica incerteza da média; diferenças pareadas isolam o efeito da estratégia; Pearson quantifica a proxy ($r=+0{,}85$) e a independência entre cabeças; desvio da soma softmax ($10^{-16}$) calibra a implementação contra T1; latência documenta custo operacional. Cada número do Capítulo~5 deriva desses operadores aplicados aos mesmos 200 sorteios fixos.
\textit{Lemas e axiomas.} T1--T7 fixam semântica (soma, limites, máximo, invariância, pertinência, singleton); T8 justifica decidir pelo máximo; T9 cobre o top-\textit{k} operado pelo código; T10 formaliza a intuição da temperatura sem análise assintótica; T11 certifica o truque numérico em reais exatos. Os três axiomas são o mínimo padrão da biblioteca, auditados por teorema --- nenhum atalho escondido.
\textit{Softmax argumentado.} A exponencial normalizada é escolhida por três motivos demonstrados: calibração (T1/T5/T7), robustez a deslocamento (T4/T11) e monotonicidade (T8/T10); a alternativa linear careceria dos três.
\textit{Comparações.} Calibração: a literatura aprende temperaturas e relaciona perdas próprias \cite{komisarenko2024,xuan2025,xie2024}, enquanto aqui a temperatura é teorema elementar sem aprendizagem --- complementar, não concorrente. Floats: aproximação correta e certificação via Lipschitz \cite{bagnara2019,murray2026} miram garantias que este trabalho não reivindica; contribuímos o lema exato e o fuzz. Orquestração: benchmarks multidimensionais e lacunas de perspectiva \cite{agentbench2023,perspectivegap2026}, recompensa de orquestração \cite{orchrm2026} e coordenação atencional \cite{orchestrator2025} definem o palco externo onde este roteador ainda não pisou. (Síntese por resumos.)
\textit{Estado da arte deste trabalho.} Numa matriz formal $\times$ empírico $\times$ auditável: formal coberto para o núcleo (11/11 verificados); empírico coberto internamente (bancada, ablações, fuzz, correlações); auditável coberto (axiomas, alegações, referências, pins, sincronia); externo descoberto (sem head-to-head, sem reprodução independente, sem Q/K/V formal completo). A novidade reivindicável é composicional --- verificação, calibração e auditoria amarradas sobre um mesmo objeto com trilha pública --- não a de cada peça isolada.
